% =====================================================================
% Report template: Break Through Tech AI Studio x Anote (Fall 2026)
% Section layout follows Setty et al., "Improving Retrieval for RAG based
% Question Answering Models on Financial Documents" (arXiv:2404.07221):
%   Abstract / 1 Introduction / 2 Techniques / 3 Evaluation / 4 Results /
%   5 Impact / 6 Conclusion and Next Steps / References
% Subsection titles are placeholders: rename them to fit our project.
% Build:  pdflatex main.tex  (run twice so references resolve)
% =====================================================================
\documentclass[11pt]{article}

% ---- Packages ----
\usepackage[margin=1in]{geometry}
\usepackage[T1]{fontenc}
\usepackage[utf8]{inputenc}
\usepackage{lmodern}
\usepackage{microtype}
\usepackage{amsmath,amssymb}
\usepackage{graphicx}
\usepackage{booktabs}
\usepackage{xcolor}
\usepackage{listings}
\usepackage[hidelinks]{hyperref}

\graphicspath{{figures/}}

% Red placeholder text; search for "TODO" to find what is left to write
\newcommand{\todo}[1]{\textcolor{red}{[TODO: #1]}}

% For prompts and code snippets
\lstset{basicstyle=\ttfamily\small, breaklines=true, frame=single,
        columns=fullflexible}

% ---- Title block ----
\title{\todo{Paper Title}}
\author{
  Luke Wang \and Brandon Romero \and Sabad Modi \and Sophia Zhang \and
  Angela Liu \and Natan Vidra
}
\date{}   % leave empty for no date, or \today

\begin{document}
\maketitle

% =====================================================================
\begin{abstract}
\todo{150--250 words: the problem, what we did, the main result (with numbers),
and why it matters.}
\end{abstract}

% =====================================================================
\section{Introduction}
\label{sec:intro}

\todo{Motivate the problem: why evaluating and improving AI coding assistants
matters, and what is hard about it. End with a short list of our contributions.}

\subsection{Background Topic 1}
\label{sec:intro-bg1}
\todo{e.g.\ Large language models for code generation and repair.}

\subsection{Background Topic 2}
\label{sec:intro-bg2}
\todo{e.g.\ Prompt optimization / evolutionary prompt search.}

% Example figure: an overview of the pipeline
\begin{figure}[ht]
  \centering
  % Replace the box with: \includegraphics[width=0.8\linewidth]{pipeline.pdf}
  \fbox{\parbox[c][4cm][c]{0.8\linewidth}{\centering Figure placeholder}}
  \caption{\todo{Overview of the pipeline.}}
  \label{fig:pipeline}
\end{figure}

\subsection{Limitations of Current Approaches}
\label{sec:intro-limits}
\todo{What existing methods or benchmarks miss, which motivates our work.}

% =====================================================================
\section{Techniques}
\label{sec:techniques}

\todo{Introduce the techniques we built or compared. Rename this section if a
more specific title fits (the source paper uses ``Techniques to Improve
Retrieval'').}

\subsection{Technique 1}
\label{sec:tech1}
\todo{Describe it. Cite related work like this~\cite{chen2021codex}.}

\subsection{Technique 2}
\label{sec:tech2}
\todo{Describe it.}

\subsection{Technique 3}
\label{sec:tech3}
\todo{Describe it.}

\subsection{Technique 4}
\label{sec:tech4}
\todo{Describe it. Delete unused subsections.}

% Example of a prompt shown verbatim
\begin{lstlisting}[caption={\todo{Example prompt.}},label={lst:prompt}]
You are an expert Python programmer. Fix the bug in the function below.
\end{lstlisting}

% =====================================================================
\section{Evaluation}
\label{sec:evaluation}

\todo{Dataset(s), models, and how we measure performance.}

\subsection{Metric 1}
\label{sec:eval-metric1}
\todo{e.g.\ functional correctness, Pass@1~\cite{chen2021codex}:}
\begin{equation}
  \text{pass@}k = \mathbb{E}\left[1 - \frac{\binom{n-c}{k}}{\binom{n}{k}}\right]
  \label{eq:passk}
\end{equation}

\subsection{Metric 2}
\label{sec:eval-metric2}
\todo{e.g.\ structural similarity such as CodeBLEU~\cite{ren2020codebleu}.}

% =====================================================================
\section{Results}
\label{sec:results}

\todo{Main findings. Refer to the figure and table:
Figure~\ref{fig:results} and Table~\ref{tab:results}.}

\begin{figure}[ht]
  \centering
  % \includegraphics[width=0.8\linewidth]{results.pdf}
  \fbox{\parbox[c][5cm][c]{0.8\linewidth}{\centering Results plot placeholder}}
  \caption{\todo{Aggregate results across methods.}}
  \label{fig:results}
\end{figure}

\begin{table}[ht]
  \centering
  \caption{\todo{Aggregate results (mean $\pm$ standard deviation).}}
  \label{tab:results}
  \begin{tabular}{lccc}
    \toprule
    Method     & Metric 1        & Metric 2        & Metric 3        \\
    \midrule
    Baseline   & 0.00 $\pm$ 0.00 & 0.00 $\pm$ 0.00 & 0.00 $\pm$ 0.00 \\
    Method A   & 0.00 $\pm$ 0.00 & 0.00 $\pm$ 0.00 & 0.00 $\pm$ 0.00 \\
    Method B   & 0.00 $\pm$ 0.00 & 0.00 $\pm$ 0.00 & 0.00 $\pm$ 0.00 \\
    \bottomrule
  \end{tabular}
\end{table}

% =====================================================================
\section{Impact}
\label{sec:impact}

\todo{Who benefits and how: e.g.\ what this means for Anote's coding
assistant and for developers using it.}

% =====================================================================
\section{Conclusion and Next Steps}
\label{sec:conclusion}

\todo{Summarize the findings in a few sentences, then list next steps.}

% =====================================================================
% References. Numbered, in order of first citation. Add entries as needed.
% (To switch to a .bib file: replace this block with
%  \bibliographystyle{unsrt} and \bibliography{references}.)
\begin{thebibliography}{9}

\bibitem{chen2021codex}
M.~Chen et~al.
\newblock Evaluating large language models trained on code.
\newblock \emph{arXiv preprint arXiv:2107.03374}, 2021.

\bibitem{ren2020codebleu}
S.~Ren et~al.
\newblock CodeBLEU: a method for automatic evaluation of code synthesis.
\newblock \emph{arXiv preprint arXiv:2009.10297}, 2020.

\end{thebibliography}

\end{document}
